\documentclass[12pt]{article}
\usepackage{amsfonts,amsthm,amsmath,amssymb,mathtools}
\usepackage{mathrsfs}
\usepackage{enumitem}
\usepackage{tikz-cd}
\usepackage{geometry}
\usepackage{mlmodern}
\usepackage{xcolor}
\usepackage{pgfplots}

\geometry{letterpaper, portrait, margin=1in}
\pgfplotsset{compat=1.18}

\setcounter{section}{0}

\renewcommand{\thesection}{\S\arabic{section}}
\renewcommand{\thesubsection}{\arabic{section}}

\newtheorem{thm}{Theorem}
\newenvironment{theorem}[1]{
	\renewcommand{\thethm}{#1}
	\begin{thm}
		}{\end{thm}}

\newtheorem{lma}{Lemma}
\newenvironment{lemma}[1]{
	\renewcommand{\thelma}{#1}
	\begin{lma}
		}{\end{lma}}

\theoremstyle{definition}
\newtheorem{ex}{}
\newenvironment{exercise}[1]{
	\renewcommand{\theex}{#1}
	\begin{ex}
		}{\end{ex}}

\title{Homework Assignment 3}
\author{Connor Mooneyhan}
\date{September 12, 2025}

\begin{document}

\maketitle

\begin{ex}
	(1.B.5)
	Give an example of a sequence of continuous real-valued functions $f_1, f_2, \dots$ on $[0, 1]$ and a continuous real-valued function $f$ on $[0,1]$ such that \begin{equation*}
		f(x) = \lim_{k \to \infty} f_k(x)
	\end{equation*}
	for each $x \in [0, 1]$ but \begin{equation*}
		\int_0^1 f \neq \lim_{k \to \infty} \int_0^1 f_k.
	\end{equation*}
\end{ex}
\begin{proof}
	Define \begin{equation*}
		f_n(x) = \begin{cases}
			2n^2x       & 0 \leq x < \frac{1}{2n}           \\
			-2n^2x + 2n & \frac{1}{2n} \leq x < \frac{1}{n} \\
			0           & \frac{1}{n} \leq x \leq 1,
		\end{cases}
	\end{equation*}
	so that $\int_0^1 f_n = \frac{1}{2}$ for all $n \in \mathbb{Z}^+$ as shown below.
	\begin{center}
		\begin{tikzpicture}
			\begin{axis}[
					axis y line=middle,
					axis x line=bottom,
					xmin=0, xmax=1,
					ymin=0, ymax=2,
					xtick={0, 1/4, 2/4, 1},
					xticklabels={$0$, $\frac{1}{2n}$, $\frac{1}{n}$, 1},
					ytick={2},
					yticklabels={$n$},
				]
				\addplot[domain=0:{1/4}, smooth, thick]{8*x};
				\addplot[domain={1/4}:{1/2}, smooth, thick]{-8*x + 4};
				\addplot[domain={1/2}:1, smooth, thick]{0};
			\end{axis}
		\end{tikzpicture}
	\end{center}
	But the limit $f_n$ as $n \to \infty$ is $f(x) = 0$, since $f_n(0) = f_n(1) = 0$ for all $n$, and for any $x \in (0, 1)$ there is an $n$ large enough that $\frac{1}{n} < x$, so $f_m(x) = 0$ for all $m \geq n$.
	Since $\int_0^1 f = 0$ and $\lim\limits_{n \to \infty} \int_0^1 f_n = \frac{1}{2}$, we have our contradiction.
\end{proof}

\begin{ex}
	(2.A.9)
	Prove that $|A| = \lim\limits_{t \to \infty}|A \cap (-t, t)|$ for all $A \subset \mathbb{R}$.
\end{ex}
\begin{proof}
	\textcolor{red}{UNFINISHED}

\end{proof}

\begin{ex}
	(2.A.11)
	Prove that if $I_1, I_2, \dots $ is a disjoint sequence of open intervals, then \begin{equation*}
		\left| \bigcup_{k=1}^\infty I_k \right| = \sum_{k=1}^\infty \ell(I_k).
	\end{equation*}
\end{ex}
\begin{proof}
	\textcolor{red}{UNFINISHED}
	This follows immediately upon noticing two key relationships.
	First, by countable subadditivity, \begin{equation*}
		\left| \bigcup_{k=1}^\infty I_k \right| \leq \sum_{k=1}^\infty \ell(I_k).
	\end{equation*}
	Second, notice that the collection of intervals $I = \left\lbrace I_k \right\rbrace_{k \in \mathbb{Z}^+}$ is itself an open cover of $\cup_{k=1}^\infty I_k$.
	By definition, for any open cover $J = \left\lbrace J_k \right\rbrace_{k \in \mathbb{Z}^+}$ of $\bigcup_{k=1}^\infty I_k$ consisting of open intervals, define an open cover $J' = \left\lbrace J_k \cap \bigcup I \right\rbrace$, which still consists of open intervals.
\end{proof}

\begin{ex}
	(2.B.5)
	Suppose $S$ is the smallest $\sigma$-algebra on $\mathbb{R}$ containing $\left\lbrace (r, r+1) : r \in \mathbb{Q} \right\rbrace$.
	Prove that $S$ is the collection of Borel subsets of $\mathbb{R}$.
\end{ex}
\begin{proof}
	By the definition Axler gives for Borel sets, we want to prove that the smallest $\sigma$-algebra on $\mathbb{R}$ containing $\mathcal{R} = \left\lbrace (r, r+1) : r \in \mathbb{Q} \right\rbrace$ is in fact the smallest $\sigma$-algebra containing all open subsets of $\mathbb{R}$.
	Let $\mathcal{B}$ be the collection of Borel sets.
	Since the elements of $\mathcal{R}$ are open, certainly $\mathcal{R} \subset \mathcal{B}$.

	On the other hand, consider the fact that finite intersections of elements of $\mathcal{R}$ are also in $\mathcal{R}$ to see that $(p, q) \in \mathcal{R}$ where $p, q \in \mathbb{Q}$ such that $q - p < 1$, since $(p, q) = (p, p + 1) \cap (q - 1, q)$.
	The set of all such ``rational intervals'', call it $\mathscr{U}$, is a basis for the topology of $\mathbb{R}$ (the proof of this fact is routine, but something I'm assuming we can brush past for the purposes of this course).
	The important thing to realize here is that the open sets of $\mathbb{R}$ can be expressed as unions of elements of $\mathscr{U}$, and since $\mathscr{U}$ is countable, the unions of elements of $\mathscr{U}$ can be expressed as countable unions.
	Thus, every open set is a countable union of elements of $\mathscr{U} \subset \mathcal{R}$, so every open set is in $\mathcal{R}$.
	This makes $\mathcal{R}$ a $\sigma$-algebra containing all open sets, so $\mathcal{B}$, which is the \textit{smallest} $\sigma$-algebra containing all open sets, is a subset of $\mathcal{R}$.
\end{proof}

\begin{ex}
	(2.B.11)
	Suppose $\mathcal{T}$ is a $\sigma$-algebra on a set $Y$ and $X \in \mathcal{T}$.
	Let $S = \left\lbrace E \in \mathcal{T} : E \subset X \right\rbrace$. \begin{enumerate}[label=(\alph*)]
		\item Show that $S = \left\lbrace F \cap X : F \in \mathcal{T} \right\rbrace$.
		\item Show that $S$ is a $\sigma$-algebra on $X$.
	\end{enumerate}
\end{ex}
\begin{proof}
	\begin{enumerate}[label=(\alph*)]
		\item Let $C = \left\lbrace F \cap X : F \in \mathcal{T} \right\rbrace$.
		      If $F \in \mathcal{T}$, then $F \cap X \in \mathcal{T}$ since $\sigma$-algebras are closed under finite intersections, and thus $C \subset S$.
		      On the other hand, if $E \in S$, then $E \in \mathcal{T}$ and $E \subset X$, so $E = E \cap X$, and we conclude that $S \subset C$.
		\item Let $E \in \mathcal{T}$, $E \subset X$.
		      See that \begin{align*}
            X \setminus E & = X \cap (Y \setminus E),
		      \end{align*}
          and since both $X$ and $Y \setminus E$ are in $\mathcal{T}$, $X \setminus E \in \mathcal{T}$.
          Also, $X \setminus E \subset X$, so $X \setminus E \in S$, making $S$ closed under complementation in $X$.

          Let $\left\lbrace E_i \right\rbrace_{i \in \mathbb{Z}^+} \subset \mathcal{T}$ such that for each $i$, $E_i \subset X$.
          Then $\bigcup_{i=1}^\infty E_i \in \mathcal{T}$ since $\mathcal{T}$ is closed under countable unions, and $\bigcup_{i = 1}^\infty E_i \subset X$, so $\bigcup_{i = 1}^\infty E_i \in S$.
          Thus, $S$ is a $\sigma$-algebra on $X$.
	\end{enumerate}
\end{proof}

\begin{ex}
	(2.B.13)
	Suppose $(X, S)$ is a measurable space, $E_1, \dots, E_n$ are disjoint subsets of $X$, and $c_1, \dots, c_n$ are distinct nonzero real numbers.
	Prove that $c_1 \chi_{E_1} + \cdots + c_n\chi_{E_n}$ is an $S$-measurable function if and only if $E_1, \dots, E_n \in S$.
\end{ex}
\begin{proof}
	\textcolor{red}{NOT STARTED}
\end{proof}

\end{document}
